\section{Mode D1 -- Arrêt d'urgence}
\label{secModeD1}

\subsection{Conditions d'entrée}

Le mode D1 est le mode dans lequel se trouve l'installation à la mise en
route de l'automate (premier mode actif). Il est également atteint depuis
n'importe quel autre mode dès qu'un arrêt d'urgence est enclenché ou 
qu'une défaillance de communication est
détectée avec l'un des équipements de la cellule (baie de commande du
robot, IHM ou contrôleur du convoyeur).

\subsection{Comportement attendu}

\begin{itemize}
    \item \textbf{Robot} : le bras doit être mis hors puissance, dès lors
    que les équipements nécessaires à cette opération sont eux-mêmes
    disponibles. Le système de préhension conserve l'état dans lequel il
    se trouvait.
    \item \textbf{Convoyeur / guichet} : l'accès au guichet doit être
    interdit.
    \item \textbf{IHM} : la situation d'arrêt d'urgence doit être signalée
    sans ambiguïté à l'opérateur, par une signalisation dédiée et fixe
    (non clignotante), distincte de la signalisation utilisée pour les
    autres arrêts.
\end{itemize}

\begin{UPSTIwarning}{Mode figé tant que la communication n'est pas rétablie}
L'installation reste figée dans ce mode tant que la communication avec
l'ensemble des équipements de la cellule n'est pas redevenue opérationnelle.
Aucune évolution vers un autre mode n'est possible avant ce retour à la
normale.
\end{UPSTIwarning}

\subsection{Conditions de sortie}

\begin{itemize}
    \item Vers le mode \textbf{D2} (\emph{Diagnostic / traitement de la
    défaillance}), dès que la communication avec l'ensemble des
    équipements de la cellule est de nouveau opérationnelle.
\end{itemize}
